\subsection*{Принятая модель погрешностей}
По уточнению автора погрешность отсчёта равна одной единице последнего
записанного разряда $\delta x$, приборная составляющая равна $0{,}03|x|$.
Например, для $3{,}00$ мА принимается $\delta I=0{,}01$ мА,
для $1{,}5$ мА --- $\delta I=0{,}1$ мА. Пределы трактуются как полуширины
равномерных распределений, поэтому стандартная неопределённость показания
\[
 u(x)=\sqrt{\frac{\delta x^2+(0{,}03x)^2}{3}}.
\]
Это расчётное соглашение при отсутствии паспортных допусков.
Для каждого прибора $V_1,A_1,V_2,A_2$ принято одно общее смещение шкалы,
сохраняющееся между отсчётами, полярностями и режимами. Ошибки последнего
разряда считаются независимыми. В расчёте МНК распространены ошибки обеих
координат, включая корреляции отсчёта с вычисленным началом координат.

К приборной ковариации добавлена прежняя оценка разброса относительно
регрессии. Такое сложение консервативно, поскольку остатки могут частично
включать ошибки отсчётов. Таблицы и кресты на графиках показывают
стандартные неопределённости $u$, соответствующие одному среднеквадратическому
отклонению. Все номинальные точки и интервалы регрессии сохранены.
Параметры $d=0{,}2$ мм, $l=5{,}2$ мм, давление, температура ионов
и коэффициент делителя используются как заданные номиналы установки.
Допуски этих номиналов в бюджет не включены. Изменение метода учёта
сдвига полярностей рассмотрено отдельно в обсуждении.

Напряжение зажигания с учётом принятой модели равно
$U_{\text{заж}}=$ \input{\proc{ai_ignition.tex}} кВ.
